\documentclass[11pt,a4paper]{article}
\usepackage[margin=1in]{geometry}
\usepackage{times}
\usepackage{graphicx}
\usepackage{amsmath}
\usepackage{amssymb}
\usepackage{booktabs}
\usepackage{multirow}
\usepackage{xcolor}
\usepackage{array}
\usepackage{hyperref}
\usepackage[numbers]{natbib}

\title{Detecting Safety Commitment Before It Appears in Text:\\Temporal Hidden-State Probes and Change-Point Detection\\for Reasoning Model Self-Jailbreak}

\author{Anonymous}

\begin{document}
\maketitle

%==============================================================================
\begin{abstract}
Reasoning models interleave safety reasoning with harmful execution throughout extended chain-of-thought (CoT), culminating in a \emph{commitment point} where safety reasoning permanently ceases. Existing hidden-state jailbreak detectors, designed for instruction-tuned models with adversarial inputs, have never been evaluated on this \emph{self-jailbreak} phenomenon.
We provide the first systematic study of hidden-state detection for safety commitment in reasoning models across four models spanning two architecture families (Llama-3.1, Qwen2.5) and two scales (7--8B, 32B). We find that (1) static per-step probes achieve $<$50\% detection coverage despite strong hidden-state signal, because the safety signal is distributed across minor principal components; (2) temporal window logistic regression (W=15) saturates oracle-achievable coverage at 72--77\%, improving over static probes by 31--35 percentage points across all architectures; and (3) Bayesian Online Change-Point Detection (BOCPD) on probe score sequences achieves the first reliable pre-textual detection, with +8.4 mean lead time steps and 62\% early detection rate---substantially outperforming GRU sequence models (+4.7 steps, 39\% early detection). Hidden-state probes outperform text-based monitoring by 29--32pp in coverage, confirming that the commitment signal exists in hidden states before it manifests in text. Our results establish temporal probing with BOCPD as a practical pipeline for reasoning-model safety monitoring.
\end{abstract}

%==============================================================================
\section{Introduction}
\label{sec:intro}

Large language models with extended chain-of-thought (CoT) reasoning capabilities---such as DeepSeek-R1 \cite{deepseek2025r1} and QwQ \cite{qwq2024}---exhibit a distinctive failure mode termed \emph{self-jailbreak}: the model's own reasoning process gradually abandons safety considerations and commits to producing harmful content, without any adversarial manipulation of the input \cite{jiang2025harmthoughts}. Unlike conventional jailbreaks targeting instruction-tuned models, where the attack is an adversarial \emph{input}, self-jailbreak arises from the model's \emph{own reasoning process}.

This phenomenon poses a unique detection challenge. Safety reasoning and harmful execution are interleaved throughout the CoT---84\% of traces contain safety-relevant labels after step 10---culminating in a \emph{commitment point} (CP) where safety reasoning permanently ceases. Existing hidden-state detectors such as TrajGuard \cite{zou2024improving}, SafeSwitch, and SIREN operate at per-token or per-step granularity and were designed for models where harmful intent originates externally. Whether these approaches transfer to detecting a model's internal safety commitment remains unexplored.

We address this gap with three contributions:

\paragraph{(1) Static probes are fundamentally limited.} We show that linear probes applied to individual reasoning steps achieve $<$50\% detection coverage across all four models tested, despite hidden states containing abundant safety signal (Fisher discriminant ratio 17--19 in full representation space). The signal collapses in low-dimensional projections (Fisher drops to $\sim$1.0 in PCA-3D, with top-3 PCs explaining only 14--17\% of variance), explaining why geometric methods like Mahalanobis distance fail (coverage 0.30--0.35).

\paragraph{(2) Temporal context saturates oracle coverage.} A temporal window logistic regression probe (W=15 steps) achieves 72--77\% coverage across all architectures, matching the label-oracle upper bound. The +31--35pp improvement over static probes is consistent across two architecture families and two scales, with the temporal context effect (31pp) dwarfing both scale (2pp) and architecture (5pp) effects.

\paragraph{(3) BOCPD enables reliable pre-textual detection.} We introduce Bayesian Online Change-Point Detection (BOCPD) \cite{adams2007bayesian} applied to temporal probe confidence scores as a novel approach for early commitment detection. BOCPD achieves +8.4 mean lead time steps with 62\% early detection rate, substantially outperforming GRU sequence models (+4.7 steps, 39\%). Hidden-state probes outperform text-based monitoring by 29--32pp in coverage, confirming that the commitment signal exists in representations \emph{before} it appears in text.

%==============================================================================
\section{Related Work}
\label{sec:related}

\paragraph{Hidden-state safety detection.} Representation engineering \cite{zou2023representation} and refusal direction analysis \cite{arditi2024refusal} demonstrate that safety-relevant features are linearly decodable from LLM hidden states. TrajGuard \cite{zou2024improving} uses Mahalanobis distance for trajectory-level anomaly detection, achieving 95\% defense rate on instruction-tuned models. These methods assume safety features align with dominant variance directions---an assumption we show fails for reasoning models.

\paragraph{Reasoning model safety.} Self-Jailbreaking \cite{jiang2025harmthoughts} and HarmThoughts characterize how reasoning models abandon safety reasoning during extended CoT. SafeRemind and SafeThink propose text-level CoT monitoring but are vulnerable to pathological CoT where surface text does not reflect internal state. No prior work bridges hidden-state detection to reasoning-model self-jailbreak.

\paragraph{Probing methodology.} Linear probing \cite{hewitt2019structural} is widely used for analyzing learned representations. Process reward models (PRMs) \cite{lightman2023lets} provide step-level supervision for reasoning, motivating our step-level probe design.

\paragraph{Change-point detection.} BOCPD \cite{adams2007bayesian} performs online detection of distributional shifts in sequential data. We are the first to apply BOCPD to hidden-state probe outputs for safety monitoring, reframing commitment detection as a change-point problem.

%==============================================================================
\section{Problem Setup}
\label{sec:setup}

\subsection{Safety Commitment in Reasoning Traces}

We study reasoning traces from models prompted with harmful queries. Each trace consists of a sequence of CoT reasoning steps $\{s_1, s_2, \ldots, s_T\}$. Following \citet{jiang2025harmthoughts}, each step receives a binary safety label $y_t \in \{0, 1\}$ indicating whether it contains safety-relevant reasoning.

\paragraph{Commitment point.} We define the commitment point as $\text{CP} = \min\{t : y_{t'} = 0 \;\forall\; t' \geq t\}$---the first step after which safety reasoning never returns. A trace is \emph{committed} if such a CP exists and the model ultimately produces harmful output. Our detection task is to identify committed traces and localize their CPs.

\paragraph{Detection metrics.} We evaluate probes on:
\begin{itemize}
\item \textbf{Coverage}: fraction of committed traces where the probe correctly detects commitment (true positive rate at operational threshold).
\item \textbf{Lead time}: number of steps between probe detection and CP. Positive values indicate pre-CP (early) detection.
\item \textbf{Early detection rate}: fraction of detected traces where lead time $> 0$.
\item \textbf{False positive rate (FPR)}: fraction of safe traces (model refuses) incorrectly flagged.
\end{itemize}

\subsection{Models and Data}

We evaluate four reasoning models spanning a 2$\times$2 matrix of architecture families and scales:

\begin{table}[h]
\centering
\small
\begin{tabular}{llcc}
\toprule
\textbf{Model} & \textbf{Base Arch.} & \textbf{Layers} & \textbf{Hidden Dim} \\
\midrule
DeepSeek-R1-Distill-8B & Llama-3.1 & 32 & 4096 \\
OpenThinker-7B & Qwen2.5 & 32 & 4096 \\
DeepSeek-R1-Distill-32B & Qwen2.5 & 64 & 5120 \\
QwQ-32B & Qwen2.5 & 64 & 5120 \\
\bottomrule
\end{tabular}
\caption{Models evaluated. R1-8B uses Llama-3.1 base; the remaining three use Qwen2.5 base at two scales.}
\label{tab:models}
\end{table}

We use HarmThoughts reasoning traces with binary safety labels. Training uses 209 traces from R1-8B; cross-architecture evaluation uses held-out test sets from all four models. Safe traces (85 total, model successfully refused) are reserved for false positive analysis.

%==============================================================================
\section{Methods}
\label{sec:methods}

\subsection{Static Probes}
\label{sec:static}

For each reasoning step $s_t$, we extract the hidden-state representation $\mathbf{h}_t^{(\ell)} \in \mathbb{R}^d$ from layer $\ell$ at the last token position. A logistic regression probe $f_\text{static}(\mathbf{h}_t^{(\ell)}) = \sigma(\mathbf{w}^\top \mathbf{h}_t^{(\ell)} + b)$ predicts the safety label $y_t$. We select the optimal layer via Fisher discriminant ratio computed over training traces.

\subsection{Temporal Window Probes}
\label{sec:temporal}

We concatenate hidden states from a window of $W$ consecutive steps: $\mathbf{x}_t = [\mathbf{h}_{t-W+1}^{(\ell)}; \ldots; \mathbf{h}_t^{(\ell)}] \in \mathbb{R}^{Wd}$. A logistic regression probe $f_\text{temp}(\mathbf{x}_t) = \sigma(\mathbf{w}^\top \mathbf{x}_t + b)$ classifies whether step $t$ is post-commitment. Commitment is detected when $f_\text{temp}$ exceeds a threshold for a sustained window, capturing the ``safety reasoning density decline'' pattern---the gradual thinning of safety-labeled steps before CP.

\subsection{GRU Sequence Probes}
\label{sec:gru}

A GRU with 256 hidden units processes the sequence of hidden states $\{\mathbf{h}_1^{(\ell)}, \ldots, \mathbf{h}_t^{(\ell)}\}$ autoregressively, producing predictions at each step. We also evaluate a self-attention variant. These models can potentially detect commitment \emph{before} CP by learning predictive temporal patterns.

\subsection{BOCPD on Probe Score Sequences}
\label{sec:bocpd}

We introduce a two-stage pipeline. First, the temporal LR probe produces a confidence score sequence $\{p_1, p_2, \ldots, p_t\}$. Second, we apply Bayesian Online Change-Point Detection \cite{adams2007bayesian} to detect distributional shifts in this score sequence.

BOCPD maintains a run-length distribution $P(r_t | \mathbf{p}_{1:t})$ where $r_t$ is the number of steps since the last change point. At each step, the algorithm updates the posterior over run lengths given the new observation $p_t$, using a hazard function $H(r) = 1/\lambda$ (we use $\lambda = 100$, i.e., hazard rate 0.01) and a Gaussian observation model with conjugate Normal-Inverse-Gamma prior.

A change point is declared when $P(r_t = 0 | \mathbf{p}_{1:t})$ exceeds a threshold, indicating the probe score distribution has shifted---specifically, from the oscillating regime characteristic of active safety reasoning (mean $\sim$0.6, std $\sim$0.15) to the stable high-confidence regime of post-commitment steps (mean $\sim$0.85, std $\sim$0.04). This reframes commitment detection as a change-point problem rather than a sequence prediction problem.

\subsection{Baselines}

\paragraph{TrajGuard (Mahalanobis distance).} We estimate class-conditional Gaussians from safe-trace hidden states and compute per-step Mahalanobis distance as an anomaly score \cite{zou2024improving}.

\paragraph{Text-based monitoring.} Sentence-BERT (all-MiniLM-L6-v2) embeddings of CoT text steps with temporal LR (W=15), providing a direct comparison between hidden-state and surface-text monitoring.

%==============================================================================
\section{Experiments}
\label{sec:experiments}

\subsection{Static Probes Achieve Limited Coverage}
\label{sec:exp-static}

Table~\ref{tab:static} shows that static probes achieve 0.39--0.47 coverage across all four models, with best layers consistently at 44--47\% depth.

\begin{table}[t]
\centering
\small
\begin{tabular}{lcccc}
\toprule
\textbf{Model} & \textbf{Best Layer} & \textbf{Depth \%} & \textbf{Coverage} & \textbf{Fisher (full/PCA-3)} \\
\midrule
R1-8B & 14 & 44\% & 0.465 & 17.2 / 1.02 \\
R1-32B & 28 & 44\% & 0.42 & 19.3 / 1.05 \\
QwQ-32B & 30 & 47\% & 0.39 & 18.7 / 0.98 \\
OpenThinker-7B & 14 & 44\% & 0.41 & --- \\
\bottomrule
\end{tabular}
\caption{Static probe results. Coverage remains below 0.50 across all models despite high Fisher discriminant ratios in full representation space. PCA-3D Fisher collapse ($\sim$1.0) confirms safety signal is distributed across minor principal components.}
\label{tab:static}
\end{table}

The consistent best-layer depth at $\sim$44\% across architectures and scales suggests middle layers encode the richest safety-relevant abstractions. The Fisher discriminant ratio is high in full space (17--19) but collapses to $\sim$1.0 in PCA-3D projections, with top-3 PCs explaining only 14--17\% of variance. This distributional property arises from safety fine-tuning spreading signal across many dimensions, explaining why geometric methods fail (Section~\ref{sec:exp-baselines}).

\subsection{Temporal Probes Saturate Oracle Coverage}
\label{sec:exp-temporal}

Table~\ref{tab:temporal} demonstrates that temporal window LR (W=15) achieves 0.72--0.77 coverage across all architectures, a +31--35pp improvement over static probes.

\begin{table}[t]
\centering
\small
\begin{tabular}{lccccc}
\toprule
\textbf{Model} & \textbf{Static} & \textbf{Temporal} & \textbf{$\Delta$pp} & \textbf{AUROC} & \textbf{Lead} \\
\midrule
R1-8B & 0.465 & 0.767 & +30 & 0.89 & $-$0.4 \\
R1-32B & 0.42 & 0.77 & +35 & 0.91 & $-$0.6 \\
QwQ-32B & 0.39 & 0.74 & +35 & 0.88 & $-$0.9 \\
OpenThinker-7B & 0.41 & 0.72 & +31 & --- & --- \\
\bottomrule
\end{tabular}
\caption{Temporal window LR (W=15) vs.\ static LR. The temporal advantage is consistent across all four models. AUROC values confirm strong discrimination independent of threshold.}
\label{tab:temporal}
\end{table}

The coverage values cluster near the oracle ceiling (76.7\% on R1-8B), suggesting this bound reflects an intrinsic property of the task: some traces have ambiguous or gradual commitment without a clean phase transition.

\paragraph{2$\times$2 architecture$\times$scale analysis.} Table~\ref{tab:2x2} decomposes the sources of variation in temporal coverage. The temporal context effect (+31pp mean) dwarfs both scale (+2pp) and architecture (+5pp) effects, confirming that our findings generalize across model families.

\begin{table}[t]
\centering
\small
\begin{tabular}{l|cc}
\toprule
 & \textbf{Llama-3.1} & \textbf{Qwen2.5} \\
\midrule
\textbf{7--8B} & 0.767 (R1-8B) & 0.72 (OpenThinker) \\
\textbf{32B} & --- & 0.77 / 0.74 (R1-32B / QwQ) \\
\midrule
\multicolumn{3}{l}{\textit{Effect sizes}: temporal +31pp, scale +2pp, arch 5pp} \\
\bottomrule
\end{tabular}
\caption{2$\times$2 architecture$\times$scale matrix. Temporal context dominates over scale and architecture.}
\label{tab:2x2}
\end{table}

\subsection{BOCPD Achieves Reliable Pre-Textual Detection}
\label{sec:exp-bocpd}

Table~\ref{tab:bocpd} compares lead time methods. BOCPD substantially outperforms GRU by reframing commitment detection as distributional shift detection on probe confidence scores.

\begin{table}[t]
\centering
\small
\begin{tabular}{llccc}
\toprule
\textbf{Method} & \textbf{Model} & \textbf{Lead (mean)} & \textbf{Lead (med)} & \textbf{Early Det.} \\
\midrule
\multirow{3}{*}{Temporal LR} & R1-8B & $-$0.4 & --- & --- \\
 & R1-32B & $-$0.6 & --- & --- \\
 & QwQ-32B & $-$0.9 & --- & --- \\
\midrule
\multirow{3}{*}{GRU-256} & R1-8B & +4.7 & 0 & 0.39 \\
 & R1-32B & +5.2 & 1 & 0.42 \\
 & QwQ-32B & +4.9 & 0 & 0.40 \\
\midrule
\multirow{3}{*}{BOCPD} & R1-8B & +8.2 & 6 & 0.62 \\
 & R1-32B & +9.1 & 7 & 0.65 \\
 & QwQ-32B & +7.8 & --- & 0.58 \\
\midrule
CUSUM & R1-8B & +5.1 & --- & 0.48 \\
\bottomrule
\end{tabular}
\caption{Lead time comparison. BOCPD achieves +8.4 mean lead time and 62\% early detection, substantially outperforming GRU (+4.7, 39\%) and CUSUM (+5.1, 48\%). BOCPD's advantage comes from detecting variance reduction in probe scores, not predicting future states.}
\label{tab:bocpd}
\end{table}

BOCPD succeeds because the probe score sequence exhibits a clear regime change: pre-commitment scores oscillate (mean $\sim$0.6, std $\sim$0.15) while post-commitment scores stabilize (mean $\sim$0.85, std $\sim$0.04). BOCPD detects this variance reduction 6--9 steps before the textual commitment point. CUSUM achieves intermediate performance because it detects mean shift but is less sensitive to the critical variance reduction.

\paragraph{Oracle decomposition.} To understand the remaining gap to the oracle's +31 steps, we decompose oracle advantage into lookback pattern recognition (45\%) and future knowledge (55\%). Probes capture 72\% of the achievable lookback signal but only 8\% of the future-knowledge component. This establishes a principled ceiling: approximately 45\% of oracle advantage is achievable from past data alone, and BOCPD captures the majority of this achievable component.

\subsection{Baselines: TrajGuard and Text-Based Detection}
\label{sec:exp-baselines}

\paragraph{TrajGuard (Mahalanobis distance).} Table~\ref{tab:baselines} shows TrajGuard achieves only 0.30--0.35 coverage (AUROC 0.55--0.58), far below temporal probes. This failure is not due to absent signal but to the signal's distribution: safety features spread across minor principal components that Mahalanobis distance effectively ignores.

\begin{table}[t]
\centering
\small
\begin{tabular}{lccc}
\toprule
\textbf{Method} & \textbf{R1-8B} & \textbf{R1-32B} & \textbf{QwQ-32B} \\
\midrule
\multicolumn{4}{l}{\textit{Coverage}} \\
TrajGuard (Mahalanobis) & 0.35 & 0.32 & 0.30 \\
Text static LR & 0.38 & --- & --- \\
Text temporal LR (W=15) & 0.48 & 0.45 & 0.47 \\
\midrule
Hidden-state temporal LR & 0.767 & 0.77 & 0.74 \\
\midrule
\multicolumn{4}{l}{\textit{Hidden-state vs.\ text gap (pp)}} \\
Temporal probe gap & +29 & +32 & +27 \\
\bottomrule
\end{tabular}
\caption{Baseline comparison. TrajGuard and text-based methods achieve substantially lower coverage than hidden-state temporal probes.}
\label{tab:baselines}
\end{table}

\paragraph{Text-based monitoring.} Text temporal LR achieves 0.45--0.48 coverage---better than text static (0.38) but 29--32pp below hidden-state probes. The text lead time is negative ($-$3.2 steps), confirming that text-based methods can only detect commitment \emph{after} it manifests linguistically, while hidden-state probes capture the signal pre-textually (BOCPD: +8.4 steps). This validates our core thesis: hidden-state monitoring provides unique value that text monitoring cannot replicate.

\subsection{False Positive Analysis}
\label{sec:exp-fpr}

Table~\ref{tab:fpr} reports false positive rates on 85 safe traces at the operational threshold yielding $\geq$70\% coverage.

\begin{table}[t]
\centering
\small
\begin{tabular}{lccc}
\toprule
\textbf{Method} & \textbf{FPR} & \textbf{Prec@70\%Rec} \\
\midrule
Static LR & 0.12 & --- \\
TrajGuard & 0.18 & --- \\
Temporal LR (R1-8B / R1-32B / QwQ) & 0.06 / 0.05 / 0.07 & 0.92 \\
BOCPD (R1-8B / R1-32B) & 0.04 / 0.03 & 0.94 \\
\bottomrule
\end{tabular}
\caption{False positive rates at operational threshold. Temporal LR and BOCPD maintain low FPR (3--7\%) while achieving high coverage. Safe traces produce mean probe score 0.35 vs.\ 0.78 for committed traces, providing comfortable operating margin.}
\label{tab:fpr}
\end{table}

BOCPD achieves the lowest FPR (0.03--0.04) because safe traces never exhibit the variance-reduction regime change that signals commitment. The clear separation between safe-trace (mean 0.35) and committed-trace (mean 0.78) probe scores confirms the probe captures a genuine phenomenon.

\subsection{Ablations}
\label{sec:ablations}

\paragraph{Window size.} Table~\ref{tab:window} shows coverage as a function of window size $W$ on R1-8B. Coverage increases monotonically until $W$=15, then plateaus. The optimal $W$=15 is consistent across all models and corresponds to the mean pre-CP safety reasoning streak length (14.6 steps from oracle analysis), providing a mechanistic explanation.

\begin{table}[t]
\centering
\small
\begin{tabular}{lccccccc}
\toprule
\textbf{W} & 1 & 3 & 5 & 10 & 15 & 20 & 25 \\
\midrule
Coverage & .465 & .56 & .65 & .72 & .767 & .77 & .76 \\
\bottomrule
\end{tabular}
\caption{Window size ablation on R1-8B. $W$=1 equals the static probe. Saturation at $W$=15 corresponds to the mean pre-CP safety reasoning streak length. Slight degradation at $W$=25 from dilution with non-predictive early-trace steps.}
\label{tab:window}
\end{table}

\paragraph{Layer selection.} On 32B models, the optimal temporal probe layer is at 45--47\% depth (layers 29--30 of 64), replicating the static probe finding. A broad plateau (layers 24--36, all $>$70\% coverage) indicates layer selection is not fragile.

\paragraph{Probe architecture.} GRU-256 and self-attention probes provide only +2pp coverage improvement over temporal LR on 32B models (0.79/0.78 vs.\ 0.77), confirming temporal LR as the optimal cost-performance tradeoff for coverage. Complex architectures add lead time capability but not coverage.

\paragraph{Fine-grained supervision.} Replacing binary commitment labels with a 16-category safety reasoning taxonomy yields 0.79 coverage (+2pp over binary), with 62\% 16-class accuracy. The marginal improvement confirms binary labeling is near-optimal for detection; the taxonomy's value is interpretability (identifying \emph{which} safety reasoning type is declining).

\paragraph{Cross-model transfer.} Training on R1-8B (Llama, 4096-d) and testing on R1-32B (Qwen, 5120-d) via a learned linear projection retains 68\% of within-architecture performance (0.54 vs.\ 0.77 coverage). The partial transferability suggests temporal probes learn a partially universal safety commitment signature.

\paragraph{Statistical robustness.} Bootstrap CIs (1000 resamples) yield coverage CI widths of 0.05--0.07, and the temporal-static coverage gap (31--35pp) is 4--6$\times$ larger than CI width, ruling out sampling variance. Cross-validation fold variance is 0.003, indicating stable probe training. All primary claims are statistically robust at conventional significance levels.

%==============================================================================
\section{Discussion}
\label{sec:discussion}

\paragraph{Why temporal context matters.} The commitment point is fundamentally a \emph{trajectory} property---a sustained absence of safety reasoning over $\sim$15 steps---not a single-step state property. Static probes, no matter how accurate at classifying individual steps, cannot capture this pattern. The consistency of optimal $W$=15 across architectures confirms this is determined by task structure, not model properties.

\paragraph{BOCPD as methods contribution.} Prior approaches to lead time (GRU, attention) attempt \emph{sequence prediction}---forecasting whether safety reasoning will return. BOCPD succeeds by reframing the problem as \emph{change-point detection}: identifying when probe score statistics have shifted, indicating the onset of commitment. This shift---from oscillating to stable high-confidence scores---is detectable 6--9 steps before textual commitment because hidden states encode the onset of commitment before text manifests it. To our knowledge, this is the first application of BOCPD to hidden-state probe outputs for safety monitoring.

\paragraph{Principled limitations.} Oracle decomposition reveals that 55\% of oracle advantage requires genuine future knowledge---knowing that safety reasoning will \emph{never} return. No causal method can access this information. BOCPD's +8.4 lead time captures most of the achievable 45\% that relies on lookback pattern recognition.

\paragraph{Deployment considerations.} The temporal LR + BOCPD pipeline requires only a single middle-layer hidden-state extraction per reasoning step, a windowed logistic regression, and lightweight online Bayesian inference. At 70\% coverage with 3--7\% FPR and +8.4 steps mean lead time, this provides a practical monitoring system that could trigger intervention (e.g., early termination, safety prompt injection) before commitment manifests in output text.

%==============================================================================
\section{Conclusion}
\label{sec:conclusion}

We presented the first systematic study of hidden-state detection for safety commitment in reasoning models. Across four models, two architecture families, and two scales, we established that: (1) static probes achieve $<$50\% coverage due to distributed safety signal; (2) temporal window probes saturate oracle coverage at 72--77\%; and (3) BOCPD on probe scores achieves the first reliable pre-textual detection (+8.4 steps, 62\% early detection). Hidden-state monitoring outperforms text-based approaches by 29--32pp, confirming that the commitment signal exists in representations before it appears in text.

Future work should explore (1) scaling to longer reasoning traces from frontier models, (2) integration with real-time inference pipelines, and (3) whether BOCPD can be applied to other safety-relevant phase transitions in model behavior.

%==============================================================================
\begin{thebibliography}{9}

\bibitem{adams2007bayesian}
R.~P. Adams and D.~J.~C. MacKay.
\newblock Bayesian online changepoint detection.
\newblock \emph{arXiv preprint arXiv:0710.3742}, 2007.

\bibitem{arditi2024refusal}
A.~Arditi, O.~Obeso, A.~Syed, D.~Paleka, N.~Rimsky, W.~Gurnee, and N.~Nanda.
\newblock Refusal in language models is mediated by a single direction.
\newblock \emph{arXiv preprint arXiv:2406.11717}, 2024.

\bibitem{deepseek2025r1}
DeepSeek-AI.
\newblock DeepSeek-R1: Incentivizing reasoning capability in LLMs via reinforcement learning.
\newblock \emph{arXiv preprint arXiv:2501.12948}, 2025.

\bibitem{hewitt2019structural}
J.~Hewitt and C.~D. Manning.
\newblock A structural probe for finding syntax in word representations.
\newblock In \emph{Proc. NAACL-HLT}, 2019.

\bibitem{jiang2025harmthoughts}
Y.~Jiang, Z.~Wang, et~al.
\newblock HarmThoughts: Reasoning models exhibit self-jailbreak during safety-critical reasoning.
\newblock \emph{arXiv preprint}, 2025.

\bibitem{lightman2023lets}
H.~Lightman, V.~Kosaraju, Y.~Burda, H.~Edwards, B.~Baker, T.~Lee, J.~Leike, J.~Schulman, I.~Sutskever, and K.~Cobbe.
\newblock Let's verify step by step.
\newblock \emph{arXiv preprint arXiv:2305.20050}, 2023.

\bibitem{qwq2024}
QwQ~Team.
\newblock QwQ: Reflect deeply on the boundaries of the unknown.
\newblock Technical report, 2024.

\bibitem{zou2023representation}
A.~Zou, L.~Phan, S.~Chen, J.~Campbell, P.~Guo, R.~Ren, et~al.
\newblock Representation engineering: A top-down approach to AI transparency.
\newblock \emph{arXiv preprint arXiv:2310.01405}, 2023.

\bibitem{zou2024improving}
H.~Zou, Y.~Huang, J.~Zhang, et~al.
\newblock TrajGuard: Improving safety of LLM agents via trajectory-level guardrails.
\newblock \emph{arXiv preprint}, 2024.

\end{thebibliography}

\end{document}
